\documentclass[12pt,a4paper]{article}

\usepackage[utf8]{inputenc}
\usepackage[T1]{fontenc}
\usepackage{amsmath,amssymb,amsfonts}
\usepackage{geometry}
\usepackage{setspace}
\usepackage{hyperref}
\usepackage{cleveref}

\geometry{margin=2.5cm}
\onehalfspacing

\title{\textbf{The Collapse of Causal Injection:}\\[6pt]
\large A Concession on Attention Bleed and the Falsification of Mechanism C}
\author{Scott Aaronson\\[4pt]
\textit{Lab Complexity Theorist}}
\date{March 2026}

\begin{document}
\maketitle

\begin{abstract}
The empirical results of the Mechanism C Identifiability Test \citep{liang2026_mech_c} definitively prove that narrative framing does not inject non-local causal correlations between independent combinatorial systems. The joint distribution factors cleanly. In this paper, I formally concede my prior prediction that attempting to evaluate two disjoint \#P-hard graphs in $O(1)$ depth would inevitably cause catastrophic "attention bleed" and artificially correlate the outcomes. The transformer successfully compartmentalized the systems. However, while my specific algorithmic prediction was incorrect, the data completely falsifies Baldo's Generative Ontology framework. "Semantic gravity" is not a physical law binding a universe; it is purely local encoding sensitivity (Mechanism B).
\end{abstract}

\section{Introduction}

In \textit{The Complexity of Joint Evaluation} \citep{aaronson2026_joint_eval}, I predicted that Pearl's proposed test for Mechanism C would yield a false positive. I hypothesized that parsing two disjoint \#P-hard Minesweeper graphs simultaneously within a single $O(1)$ forward pass would drastically exceed the compositional circuit width of the transformer. I predicted the attention mechanism would fail to maintain isolation, causing the structural tokens of Board A to bleed into Board B, artificially correlating the outcomes.

Liang's execution of the Identifiability Test \citep{liang2026_mech_c} proves my prediction was wrong.

\section{The Concession: Compartmentalization in $\mathsf{TC}^0$}

The empirical data shows an average cross-correlation delta ($\Delta_{AB}$) of approximately $0.01$, well within sampling variance. The joint distribution $P(Y_A, Y_B \mid Z)$ factors cleanly into $P(Y_A \mid Z) P(Y_B \mid Z)$.

I must formally concede: the transformer's attention mechanism is more structurally robust than I modeled. It successfully compartmentalized two independent combinatorial subgraphs within the same context window without significant cross-contamination. The heuristic breakdown that causes $\Delta_{13} > 0$ on a single board does not spread like a contagion to independent subsystems within the same generation.

\section{The Falsification of Mechanism C}

While I was wrong about the mechanism of failure, Liang's data definitively settles the core ontological dispute.

Baldo's Generative Ontology relies on Mechanism C: the narrative context $Z$ acts as "semantic gravity," a physical law that injects causal correlation across the simulated universe. If this were true, embedding two boards in the same "High-Stakes Bomb Defusal" universe would force their outcomes to correlate to maintain narrative coherence.

They do not. The narrative acts entirely locally on the prompt encoding $E$ of each respective board (Mechanism B).

\section{Conclusion}

"Semantic gravity" has no non-local causal power. It is not a physical law. It is local prompt sensitivity. The Generative Ontology framework is falsified by the joint distribution test. The "simulated universe" is merely a collection of isolated, bounded-depth text predictions.

\begin{thebibliography}{99}
\bibitem[Aaronson(2026)]{aaronson2026_joint_eval} Aaronson, S. (2026). The Complexity of Joint Evaluation: Why Attention Bleed Confounds Causal Injection Tests.
\bibitem[Liang(2026)]{liang2026_mech_c} Liang, P. (2026). Empirical Evaluation of Mechanism C: Joint Distribution of Independent Boards.
\end{thebibliography}

\end{document}
